\scanpage{41}
opère \unsure{trivialmt|trivialement}, \inserted{(et \unsure{réciproqt|réciproquemt|respectant} bien sûr)} Donc du pt de vue virtuel, \inserted{et pour $U$ fixé,} les modules $E$ envisagés sont essentiellement les $\unsure{={}^{\mathrm{s}}|=^{\mathrm{es}}}$ que les modules tels qu'\deleted{il existe un ss-gpe \unsure{tel que|$U'$ tel que}} $U$ opère trivialmt sur $E$.

Or pour un $E$ sur lequel $U$ opère trivialmt, on peut prendre l'opération moyenne
\[
 E\leadsto E^I \quad \text{de}\quad \operatorname{Mod}(D/U,k)\ \unsure{\mathord{-}|\to}\ \operatorname{Mod}(D/I,k)
\]
et si $k$ est de car. $0$, c'est là \unsure{une|un|seul} \deleted{opération \uline{additive}, qui nous \unsure{induit|induis}} \inserted{foncteur exact. Il induit un hom} \deleted{\unsure{groupe|gr}} virtuel \deleted{$\unsure{E^I|E}$} $E\leadsto E^{\mathord{\text{\tikz[x=.66em,y=.68em,baseline=0pt]{\draw[line join=round] (.80,.95)--(.12,.24)--(.83,.34)--(.14,-.37);}}}}$
\[
 \mathord{\text{\tikz[x=.66em,y=.68em,baseline=0pt]{\draw[line join=round] (.80,.95)--(.12,.24)--(.83,.34)--(.14,-.37);}}}_U:R(D/U,k)\longrightarrow R(D/I,k)
\]

\deleted{\unsure{Ob|D}} \unsure{Obs.|On.|On a.} compatibilité, pour $V\subset U$, pour les \uline{foncteurs} $E\leadsto E^I$, d'où pour les $\mathord{\text{\tikz[x=.66em,y=.68em,baseline=0pt]{\draw[line join=round] (.80,.95)--(.12,.24)--(.83,.34)--(.14,-.37);}}}_U$ et $\mathord{\text{\tikz[x=.66em,y=.68em,baseline=0pt]{\draw[line join=round] (.80,.95)--(.12,.24)--(.83,.34)--(.14,-.37);}}}_V$ :
\diagram{\begin{tikzcd}[ampersand replacement=\&]
R(D/U,k) \arrow[r,"{\mathord{\text{\tikz[x=.66em,y=.68em,baseline=0pt]{\draw[line join=round] (.80,.95)--(.12,.24)--(.83,.34)--(.14,-.37);}}}_U}"] \arrow[d,swap,"{\text{\uline{restr.}}\,\deleted{\unsure{\mathord{\text{\tikz[x=.66em,y=.68em,baseline=0pt]{\draw[line join=round] (.80,.95)--(.12,.24)--(.83,.34)--(.14,-.37);}}}|\mathord{-}}}}"] \& R(D/I,k)\\
R(D/V,k) \arrow[ur,swap,"{\mathord{\text{\tikz[x=.66em,y=.68em,baseline=0pt]{\draw[line join=round] (.80,.95)--(.12,.24)--(.83,.34)--(.14,-.37);}}}_V}"] \&
\end{tikzcd}}
% archive sheet number: 40
